% titel: Blatt 7, A1: Ableitungen
% skript: Kap. 4.1 (Ableitungsregeln, Regel 4.1.8)

% Teil I: Werkzeuge, kein Lösungsweg.

Ableiten komplexer Funktionen mit den Standard-Regeln (Produkt-, Quotienten-,
Kettenregel) sowie Kurvendiskussion von $f(x)=\ln(x^2+1)$.

\begin{werkzeugkasten}[Kettenregel]
Für Verkettungen $h(x)=u(v(x))$ gilt
\begin{equation}
  h'(x)=u'(v(x))\cdot v'(x).
\end{equation}
Wichtige Fälle: $(\ln g)'=\frac{g'}{g}$, $(\mathrm e^{g})'=g'\mathrm e^{g}$,
$(\sqrt{g})'=\frac{g'}{2\sqrt g}$, $(\sin g)'=g'\cos g$. Die äußere Ableitung
wird an $v(x)$ ausgewertet, dann mit der inneren Ableitung $v'$ multipliziert.

\textit{Skript:} Kap. 4.1, Regel 4.1.8.
\end{werkzeugkasten}

\begin{werkzeugkasten}[Kurvendiskussion von $f(x)=\ln(x^2+1)$]
Der Definitionsbereich folgt aus dem Logarithmusargument: $x^2+1>0$ gilt für
alle $x\in\mathbb R$, also ist $f$ überall definiert. Für Nullstellen löst man
$f(x)=0$, für Extrema $f'=0$ (mit Vorzeichenprüfung), für Wendepunkte $f''=0$.

\textit{Skript:} Kap. 4.2--4.3.
\end{werkzeugkasten}

\begin{achtung}[Darauf achten]
\begin{itemize}
  \item Kettenregel vollständig: äußere Ableitung \emph{an der inneren Funktion}
        plus innere Ableitung. Bei $\ln(x^2+1)$ nicht $\ln$ allein ableiten.
  \item Vorzeichen des Arguments beim $\ln$: $x^2+1$ ist stets positiv — keine
        Definitionslücken erfinden.
  \item Bei Extrema und Wendepunkten das Vorzeichen der Ableitung beidseitig
        prüfen, nicht nur das Verschwinden.
\end{itemize}
\end{achtung}
